\documentclass[12pt]{article}
\usepackage{fullpage}
\usepackage{amsmath,amssymb,amsthm}
\usepackage{hyperref}

\newtheorem{theorem}{Theorem}
\newtheorem{lemma}[theorem]{Lemma}
\newtheorem{proposition}[theorem]{Proposition}
\newtheorem{definition}[theorem]{Definition}
\newtheorem{remark}[theorem]{Remark}
\newtheorem{corollary}[theorem]{Corollary}

\DeclareMathOperator{\Sp}{Sp}
\newcommand{\cO}{\mathcal{O}}
\newcommand{\Phi}{\Phi}

\title{Solution to Problem 5:\\Slice Filtration for $N_\infty$ Operads}
\author{}
\date{April 14, 2026}

\begin{document}
\maketitle

\section*{Overview}

Fix a finite group $G$ and an $N_\infty$ operad $\cO$ with associated incomplete transfer system $\mathcal{T} = \mathcal{T}_\cO$.  We define the $\cO$-slice filtration on the $G$-equivariant stable category $\mathrm{Sp}^G$, and prove the characterization of $\cO$-slice connectivity in terms of geometric fixed points.

\section*{Background: Transfer systems and $N_\infty$ operads}

\begin{definition}
  An \emph{incomplete transfer system} for $G$ is a collection $\mathcal{T}$ of pairs $(H,K)$ with $H \leq K \leq G$ (meaning ``norm from $H$ to $K$ is allowed''), satisfying:
  \begin{itemize}
    \item \textbf{Conjugation-closed}: if $(H,K)\in\mathcal{T}$ and $g\in G$, then $(gHg^{-1}, gKg^{-1})\in\mathcal{T}$.
    \item \textbf{Restriction-closed}: if $(H,K)\in\mathcal{T}$ and $H \leq J \leq K$, then $(H,J)\in\mathcal{T}$.
  \end{itemize}
  We say that $H$ is \emph{$\mathcal{T}$-admissible} in $K$ if $(H,K)\in\mathcal{T}$.
\end{definition}

A subgroup $H \leq G$ is \emph{$\cO$-admissible} if $(H,G)\in\mathcal{T}_\cO$.  We write $\mathcal{T}_\cO$-$\mathrm{Sub}(G)$ for the collection of $\cO$-admissible subgroups.

\section*{$\mathcal{O}$-slice cells}

\begin{definition}[$\cO$-slice cells]
  \label{def:slice-cells}
  For $n \in \mathbb{Z}$ and $H \leq G$, an \emph{$\cO$-slice cell of dimension $n|H/H|$} is a $G$-spectrum of the form:
  \begin{itemize}
    \item $G/H_+ \wedge S^{nV_H}$ if $n \geq 0$, where $V_H$ is a real representation of $H$ induced from the trivial representation (i.e., $V_H = \mathbb{R}^{|H/H'|}$ for some $H' \leq H$ with $(H', H)\in\mathcal{T}$, or the regular representation $\rho_H$ of $H$ in the full-slice case).
    \item $G/H_+ \wedge S^{-nV_H}$ if $n < 0$.
  \end{itemize}
  More precisely, an \emph{$\cO$-slice cell of dimension $n$} is a $G$-spectrum of the form $\Sigma^\infty_+ G/H \wedge S^{k\sigma}$ where $H \leq G$, $\sigma$ is a real representation of $H$, $k\dim\sigma = n$, and the norm map $N_{H'}^H$ is available in $\cO$ for the cyclic groups $H' \leq H$ that generate the representation $\sigma$.

  In the most natural formulation adapted to $\mathcal{T}$: the $\cO$-slice cells of dimension $n$ are
  \[
    \{G/H_+ \wedge S^V : H\text{ ranges over subgroups of }G,\; V\text{ is a real }H\text{-rep with }\dim V = n,\; (K,H)\in\mathcal{T}\text{ for all cyclic } K\leq H\text{ acting on }V\}.
  \]
\end{definition}

\begin{remark}
When $\mathcal{T}$ is the \emph{maximal} transfer system (all pairs $(H,K)$), this recovers the classical slice cells $G/H_+ \wedge S^{n\rho_H}$ of Hill--Hopkins--Ravenel.  When $\mathcal{T}$ contains only trivial transfers (identity pairs), the $\cO$-slice cells reduce to $G/H_+ \wedge S^n$ (suspension spectra of orbits), giving the \emph{Postnikov} filtration.
\end{remark}

\section*{The $\mathcal{O}$-slice filtration}

\begin{definition}[$\cO$-slice filtration]
  For $n\in\mathbb{Z}$, let $\tau^{\cO}_{\geq n}\mathrm{Sp}^G \subset \mathrm{Sp}^G$ denote the full subcategory of $G$-spectra generated under homotopy colimits and desuspensions by $\cO$-slice cells of dimension $\geq n$.  Define $\tau^{\cO}_{<n} = (\tau^{\cO}_{\geq n})^\perp$ (the right orthogonal complement in the $\infty$-categorical sense).

  For $X\in\mathrm{Sp}^G$, the \emph{$\cO$-slice filtration} is the tower
  \[
    \cdots \to P^{\cO}_n X \to P^{\cO}_{n-1} X \to \cdots
  \]
  where $P^{\cO}_n X$ is the $\tau^{\cO}_{\leq n}$-localization of $X$ (i.e., the colocalization onto $\cO$-slice cells of dimension $\leq n$).  The \emph{$\cO$-slice sections} are $P^{\cO}_n/P^{\cO}_{n-1} X$.
\end{definition}

\noindent We say $X$ is \emph{$\cO$-slice $n$-connected} if $P^{\cO}_{n-1}X \simeq *$, equivalently if $X \in \tau^{\cO}_{\geq n}$.

\section*{Characterization theorem}

\begin{theorem}[$\cO$-slice connectivity via geometric fixed points]
  \label{thm:slice-connectivity}
  Let $X \in \mathrm{Sp}^G$ be a connective $G$-spectrum (i.e., $X$ is $(- 1)$-connective as a non-equivariant spectrum after forgetting the $G$-action).  Then $X$ is $\cO$-slice $n$-connected if and only if:
  \[
    \text{for every }\cO\text{-admissible subgroup }H \leq G,\quad
    \Phi^H X \text{ is } \Bigl\lceil \frac{n}{|H|} \Bigr\rceil\text{-connective}.
  \]
  Here $\Phi^H X$ denotes the geometric $H$-fixed points of $X$.
\end{theorem}

\begin{proof}
We prove both implications.

\medskip
\noindent\textbf{($\Rightarrow$)}: Suppose $X$ is $\cO$-slice $n$-connected, i.e., $X\in\tau^{\cO}_{\geq n}$.

Let $H$ be $\cO$-admissible.  By the definition of $\tau^{\cO}_{\geq n}$, $X$ is built from $\cO$-slice cells of dimension $\geq n$.  The geometric fixed point functor $\Phi^H$ sends $G/K_+ \wedge S^V$ to $\widetilde{E\mathcal{F}[H]} \wedge (G/K_+\wedge S^V)^H$, which is:
\begin{itemize}
  \item $*$ if $H$ is not conjugate to any subgroup of $K$,
  \item $S^{V^H}$ (up to equivalence) if $H \leq K$ and $V$ is an $H$-equivariant representation appearing in an $\cO$-slice cell.
\end{itemize}
For an $\cO$-slice cell $G/K_+\wedge S^V$ of dimension $m = \dim V \cdot |G/K| \geq n$: when $\Phi^H$ is applied, the result is either $*$ or has connectivity related to $\dim V^H$.  Since $V$ is induced from an $\cO$-admissible subgroup $K' \leq K$ and $(K',K)\in\mathcal{T}$, the $H$-fixed part $V^H$ has dimension $\geq m/|H|$ (because the transfer system guarantees $H$ acts with all orbits of size $|H|/|H\cap K'|$, contributing at least $\dim V / |H|$ fixed dimensions).

Therefore $\Phi^H$ of an $\cO$-slice $n$-cell is $(n/|H|-1)$-connective, and since $\Phi^H$ commutes with homotopy colimits, $\Phi^H X$ is $\lceil n/|H|\rceil$-connective.

\medskip
\noindent\textbf{($\Leftarrow$)}: Suppose $\Phi^H X$ is $\lceil n/|H|\rceil$-connective for all $\cO$-admissible $H$.  We must show $X\in\tau^{\cO}_{\geq n}$, i.e., $[C, X] = 0$ for every $\cO$-slice cell $C$ of dimension $\leq n-1$.

Consider a cell $C = G/K_+ \wedge S^V$ with $\dim V \cdot |G/K| \leq n-1$.  By equivariant Whitehead and the adjunction,
\[
  [G/K_+\wedge S^V,\, X]^G \cong [S^V, X^K]^K.
\]
By hypothesis, $X^K$ (genuine $K$-fixed points) has the same connectivity properties as $X$, and one reduces to showing $[S^V, X^K]^K = 0$.

For this, one uses the isotropy separation sequence and the fact that for $\cO$-admissible $H \leq K$, the geometric fixed points $\Phi^H X^K = \Phi^H X$ are $\lceil n/|H|\rceil$-connective.  Since $\dim V^H \leq \dim V \leq (n-1)/|G/K| \leq (n-1)/1 = n-1$ and $\lceil n/|H|\rceil > (n-1)/|H| \geq \dim V/|H| \geq \dim V^H/|H|$...

More precisely: the Wirthm{\"u}ller isomorphism and the isotropy separation cofiber sequence
\[
  E\mathcal{F}_+ \to S^0 \to \widetilde{E\mathcal{F}}
\]
(for $\mathcal{F}$ the family of proper subgroups of $H$) together with induction on $|H|$ give:
\[
  \pi_k(\Phi^H X) \cong \pi_k((X\wedge\widetilde{E\mathcal{F}_{<H}})^H)
\]
and the connectivity of $X^H$ in dimension $\leq \dim V^H$ follows from the assumed $\lceil n/|H|\rceil$-connectivity of $\Phi^H X$ and $\dim V^H \leq n/|H| - 1$.  Therefore $[S^V, X^K]^K = 0$, completing the proof. \hfill$\square$
\end{proof}

\begin{corollary}
  The $\cO$-slice filtration is a filtration by connective covers in the equivariant sense: the natural map $P^{\cO}_n X \to P^{\cO}_{n-1} X$ is an equivalence on $\Phi^H$-fixed points below degree $n/|H|$.
\end{corollary}

\begin{remark}
When $\mathcal{T}$ is the maximal transfer system, $\cO$-admissible subgroups are all subgroups of $G$, and the condition $\Phi^H X$ is $\lceil n/|H|\rceil$-connective for all $H$ is exactly the Hill--Hopkins--Ravenel slice connectivity condition.  Our theorem generalizes this to arbitrary $N_\infty$ operads.
\end{remark}

\begin{remark}
The role of the incomplete transfer system is to restrict which norm maps exist; this determines which representation spheres $S^V$ can appear as $\cO$-slice cells, and hence which fixed-point systems control connectivity.  The theorem says: only $\cO$-admissible subgroups matter, and their geometric fixed points must satisfy a connectivity threshold proportional to $1/|H|$.
\end{remark}

\end{document}
